\section{\textcolor{red}{Экспериментальное исследование}}
\label{sec:experiments}

В предыдущих главах был формально определён \textcolor{red}{орбитальный оцениватель}
долгосрочного эффекта на основе метрик $r_V$~\eqref{eq:rel_diff} и
$r_{V\text{-corr}}$~\eqref{eq:rel_diff_corrected}, описана процедура
построения орбиталей (раздел~\ref{sec:segmentation}) и сформулирована
гипотеза о компромиссе между смещением и дисперсией
(раздел~\ref{sec:estimation_tasks}). Настоящая глава замыкает эту
конструкцию эмпирически.

Исследование преследует четыре последовательные цели:
\begin{enumerate}
    \item эмпирически проверить мультипликативное допущение
          $V_c(t) = V_c \cdot V_o(t)$ из раздела~\ref{subsec:lltv};
    \item оценить точность орбитальных прогнозов $r_V$ и
          $r_{V\text{-corr}}$ относительно фактического lLTV‑эффекта,
          наблюдаемого на горизонте $\mathcal{F}(T_E, T_F)$;
    \item проверить согласованность орбитальных оценок с
          краткосрочными метриками, уже применяемыми в продуктовой
          практике;
    \item квантифицировать компромисс смещение–дисперсия путём
          сопоставления с двумя эталонными подходами, введёнными в
          разделе~\ref{sec:estimation_tasks}.
\end{enumerate}

Все эксперименты проводятся на крупной платформе электронной
коммерции; объём каждого эксперимента — миллионы пользователей. Если
не оговорено иное, прогнозный горизонт $T_F = 30$ дней.

\textbf{Отсутствие утечки информации.} Данные на прогнозном горизонте
$\mathcal{F}(T_E, 30)$ не используются при построении орбитальной
оценки. Метод опирается только на признаки пользователей до начала
эксперимента, их орбитальные метки в моменты $T_S$ и $T_E$, а также
на наблюдаемые переходы внутри экспериментального окна $[T_S,T_E]$.
Реализованная ограниченная пожизненная ценность (lLTV) на горизонте
$\mathcal{F}(T_E, 30)$ используется исключительно как отложенный
эталон для валидации.


\subsection{Корпус экспериментов и схема оценивания}
\label{subsec:corpus_n_setup}

Эмпирическая база — корпус значимых A/B‑экспериментов, для которых
доступна верифицированная реализованная lLTV $R_u(T_E)$ на горизонте
$\mathcal{F}(T_E, 30)$. Различные подмножества корпуса используются
для разных проверок: для сопоставления предсказанного и фактического
эффекта (раздел~\ref{subsec:results_prediction}) отбирается $21$
эксперимент, удовлетворяющий условиям полноты данных; для проверки
согласованности с краткосрочной метрикой — $28$ экспериментов.

Рассматриваемые интервенции охватывают качественно разные продуктовые
воздействия: ценообразование, оптимизацию интерфейса и пользовательских
сценариев, обновления моделей поиска и ранжирования, а также запуск
новых функциональностей — в частности, опции «Buy Now, Pay Later»
(BNPL), подробно разбираемой в~\ref{subsec:bnpl_case}.

Для каждого эксперимента доступны:
\begin{itemize}
    \item поведенческие признаки пользователей за 30 дней до $T_S$,
          в соответствии с разделом~\ref{subsubsec:feature_engineering}:
          \texttt{recency}, \texttt{frequency}, \texttt{GMV} и другие
          агрегированные характеристики поведения;
    \item орбитальные метки $c_u(T_S)$ и $c_u(T_E)$, полученные по
          сегментации, описанной в разделе~\ref{sec:segmentation};
          в качестве сегментов используются
          \texttt{weak}, \texttt{active}, \texttt{paid},
          \texttt{active\&paid}, \texttt{whale} и вспомогательный
          \texttt{out};
    \item эмпирические матрицы переходов
          $\widehat{P}^{(g)}_{c\to c'}$ для $g\in\{A,B\}$ и
          производные распределения $D_g(T_S)$, $D_g(T_E)$;
    \item реализованная lLTV $R_u(T_E)$ для пользователей обеих
          групп, используемая исключительно как отложенный эталон.
\end{itemize}

В соответствии с разделом~\ref{sec:estimation_tasks} орбитальные
множители $\{\hat{V}_c\}_{c\in\mathcal{C}}$ оцениваются на историческом
окне, не пересекающемся с экспериментальным; их устойчивость во
времени проверяется в~\ref{subsec:assumption_check}.

Сопоставление ведётся с двумя эталонами:
\begin{description}
    \item[Прямое наблюдение lLTV:] относительная разница средних
          $R_u(T_E)$ между группами
          (раздел~\ref{sec:estimation_tasks}, пункт~(1));
    \item[Краткосрочная суррогатная метрика:] средняя выручка на
          пользователя в окне $[T_S, T_E]$ (\textit{Average GMV per
          User}); используется для проверки согласованности с
          действующими продуктовыми показателями.
\end{description}

Все статистические выводы получаются непараметрической
бутстрэп‑процедурой с $10\,000$ псевдовыборок. На рисунках приведены
70\% бутстрэп‑интервалы — они служат для визуального сопоставления
точечных оценок; решения о статистической значимости далее принимаются
отдельно при уровне значимости $\alpha = 0.1$.


\subsection{Эмпирическая проверка мультипликативного допущения}
\label{subsec:assumption_check}

Допущение $V_c(t) = V_c \cdot V_o(t)$, введённое в
разделе~\ref{subsec:lltv}, играет роль ключевого структурного
предположения: именно оно обеспечивает сокращение общего временного
множителя $V_o(t)$ в формуле~\eqref{eq:rel_diff} и позволяет
переводить изменения распределения по орбиталям в оценку
относительного lLTV‑эффекта. Поскольку допущение является модельным,
его эмпирическая проверка необходима до содержательного использования
$\hat{V}_c$.

Для проверки используется историческое окно длиной 1.5 года, не
пересекающееся с экспериментальным корпусом. Для каждой орбитали
$c\in\mathcal{C}$ и каждого опорного момента $t$ вычисляется оценка
$\hat{V}_c(t)$, после чего анализируются попарные отношения
$\hat{V}_c(t)/\hat{V}_{c'}(t)$ во времени
(см.~рис.~\ref{fig:monthly_tltv_dynamic}).

Результаты показывают, что относительный порядок орбиталей по lLTV и
масштаб различий между ними остаются устойчивыми, несмотря на
сезонные колебания и системные изменения платформы. Это позволяет:
\begin{itemize}
    \item принять мультипликативное допущение как рабочее
          приближение в рассматриваемом домене;
    \item использовать единый набор $\{\hat{V}_c\}_{c\in\mathcal{C}}$,
          агрегирующий $\hat{V}_c(t)$ по разным $t$, при подстановке
          в $r_V$ и $r_{V\text{-corr}}$.
\end{itemize}

Тем самым эмпирически подтверждается и более общий тезис из
вводной части раздела~\ref{sec:framework}: долгосрочное воздействие
проявляется в первую очередь через перераспределение пользователей
\textit{между} орбиталями, а не через изменение поведения внутри них.


\subsection{Точность и надёжность долгосрочных прогнозов}
\label{subsec:results_prediction}

Под \textit{lLTV‑эффектом} здесь и далее понимается относительная
разница ожиданий lLTV между тестовой и контрольной группами на
горизонте $\mathcal{F}(T_E, 30)$ — в обозначениях
раздела~\ref{subsec:lltv} это $r_V$ или $r_{V\text{-corr}}$ в
зависимости от формулировки.


\paragraph{Сопоставление предсказанного и фактического эффекта.}

В качестве отправной точки рассматривается revenue‑based вариант
методики, в котором роль целевой величины при построении орбиталей
(в смысле раздела~\ref{subsubsec:target_definition}) играет
накопленная ограниченная LTV. На рис.~\ref{fig:gmv_val} для $21$
эксперимента сопоставлены предсказанный $\hat{r}_{V\text{-corr}}$ и
фактический lLTV‑эффект.

\begin{comment}
\begin{figure}[t]
    \centering
    \begin{subfigure}[t]{0.48\linewidth}
        \includegraphics[width=\linewidth]{data/new/05_gmv_val_v4.png}
        \caption{revenue-based}
        \label{fig:gmv_val}
    \end{subfigure}
    \hfill
    \begin{subfigure}[t]{0.48\linewidth}
         \includegraphics[width=\linewidth]{data/new/05_frequency_val_v4.png}
         \caption{frequency-based}
         \label{fig:frequency_val}
    \end{subfigure}    
    \caption{Validation of Orbital-based predictions against actual lLTV uplift.
    (a) Uplift in revenue; (b) Uplift in active days.
    Horizontal and vertical lines indicate 70\% \textit{bootstrap}
    confidence intervals. The dashed lines denote \( y = x,\, y = 0,\, x = 0 \).}
    \label{fig:val}
\end{figure}
\end{comment}

В рассматриваемой выборке знаки предсказанного и фактического эффектов
в большинстве экспериментов совпадают. Вместе с тем фактические оценки
по $R_u(T_E)$ обладают высокой дисперсией: бутстрэп‑интервалы
наблюдаемого GMV‑эффекта нередко оказываются широкими, что
соответствует характеристике прямого эталона из
раздела~\ref{sec:estimation_tasks}, пункт~(1) (тяжёлые хвосты
распределения выручки), и затрудняет однозначную интерпретацию
экспериментов с умеренным эффектом.

Для дополнительной проверки направления и относительной силы эффекта
используется frequency‑based вариант методики, в котором целевой
величиной выступает число активных дней — вторая иллюстративная
целевая величина из раздела~\ref{subsubsec:target_definition}. Эта
величина не заменяет финансовую lLTV, но обладает существенно меньшей
дисперсией и даёт независимую проверку поведенческой составляющей
эффекта в случаях, когда GMV‑оценка слишком волатильна
(см.~рис.~\ref{fig:frequency_val}).

Согласованность статистических выводов по знаку и значимости
дополнительно иллюстрируется матрицами ошибок
(табл.~\ref{tbl:conf_matrix}) при уровне значимости $\alpha = 0.1$.

\begin{comment}
\begin{table*}[htbp]
    \centering
    \begin{minipage}{0.48\textwidth}
        \centering
        \scalebox{0.8}{\begin{NiceTabular}{c c c c c}
            \toprule
            & & \multicolumn{3}{c}{Prediction lLTV Uplift} \\
            \cmidrule{3-5}
            & & stat sig positive & not stat sig & stat sig negative \\
            \midrule
            \Block{3-1}{\rotate \makebox[0pt]{\shortstack{Actual \\ lLTV \\ Uplift}}}
                & stat sig positive & \cellcolor{cwred}4.76\%   & \cellcolor{cwnormal}0.00\% & \cellcolor{cwnormal}0.00\% \\
                & not stat sig      & \cellcolor{cwlightred}9.52\% & \cellcolor{cwred}66.67\% & \cellcolor{cwlightred}14.29\% \\
                & stat sig negative & \cellcolor{cwnormal}0.00\% & \cellcolor{cwnormal}0.00\% & \cellcolor{cwred}4.76\% \\
            \bottomrule
        \end{NiceTabular}}
        \vspace{1em}
        \textbf{(a) revenue-based}
    \end{minipage}
    \begin{minipage}{0.48\textwidth}
        \centering
        \scalebox{0.8}{\begin{NiceTabular}{c c c c c}
            \toprule
            & & \multicolumn{3}{c}{Prediction lLTV Uplift} \\
            \cmidrule{3-5}
            & & stat sig positive & not stat sig & stat sig negative \\
            \midrule
            \Block{3-1}{\rotate \makebox[0pt]{\shortstack{Actual \\ lLTV \\ Uplift}}}
                & stat sig positive & \cellcolor{cwlightblue}9.52\% & \cellcolor{cwlightblue}4.76\% & \cellcolor{cwnormal}0.00\% \\
                & not stat sig      & \cellcolor{cwblue}23.81\%    & \cellcolor{cwblue}38.10\% & \cellcolor{cwlightblue}4.76\% \\
                & stat sig negative & \cellcolor{cwnormal}0.00\% & \cellcolor{cwlightblue}9.52\% & \cellcolor{cwlightblue}9.52\% \\
            \bottomrule
        \end{NiceTabular}}
        \vspace{1em}
        \textbf{(b) frequency-based}
    \end{minipage}
    \caption{
        Confusion matrix of the statistical conclusions (at significance level $\alpha = 0.1$)
        drawn using the observed (vertically) and the predicted (horizontally) lLTV Uplift
        for 21 high-impact A/B tests.
    }
    \label{tbl:conf_matrix}
\end{table*}
\end{comment}

В revenue‑based варианте все статистически значимые положительные
прогнозы соответствуют неотрицательным фактическим эффектам:
противоположных по знаку статистически значимых решений в
рассматриваемом корпусе не наблюдается. Это согласуется с тем,
что было заявлено в~\ref{sec:estimation_tasks} применительно к
снижению риска ошибочных рекомендаций.

Frequency‑based вариант обладает более высокой чувствительностью:
он выявляет больше положительных эффектов ценой нескольких
расхождений с наблюдаемыми статистическими выводами по выручке;
противоположных по знаку статистически значимых решений он также не
порождает.

В совокупности на рассматриваемом корпусе орбитальный подход
демонстрирует три свойства, обещанных в~\ref{sec:estimation_tasks}:
\begin{enumerate}
    \item \textbf{Согласованность по направлению:} ни одно
          статистически значимое положительное предсказание не
          соответствует статистически значимому отрицательному
          фактическому эффекту;
    \item \textbf{Снижение дисперсии:} 70\% бутстрэп‑интервалы
          $r_{V\text{-corr}}$ оказываются, как правило, в
          $\approx 30$--$50\%$ уже соответствующих интервалов прямого
          эталона по $R_u(T_E)$;
    \item \textbf{Повышенная чувствительность:} анализ переходов
          между орбиталями выявляет устойчивые сдвиги, не различимые
          по краткосрочным агрегатам GMV.
\end{enumerate}


\paragraph{Согласованность с краткосрочной суррогатной метрикой.}

Орбитальные оценки сопоставляются со средней выручкой на пользователя
в окне $[T_S, T_E]$ на $28$ экспериментах
(см.~табл.~\ref{tbl:sign_validation_avg_gmv}).

\begin{comment}
\begin{table}[ht]
\centering
\scalebox{0.8}{\begin{NiceTabular}{c c c c c}
\toprule
 & & \multicolumn{3}{c}{Prediction lLTV Uplift} \\
\cmidrule{3-5}
 & & stat sig pos & not stat sig & stat sig neg \\
\midrule
\Block{3-1}{\rotate \makebox[0pt]{ \shortstack{Average \\ GMV \\ per user}}}
    & stat sig positive & \cellcolor{cwlightred}10.71\% & \cellcolor{cwlightred}10.71\% & \cellcolor{cwnormal}0.00\% \\
    & not stat sig      & \cellcolor{cwred}21.43\%      & \cellcolor{cwred}42.86\% & \cellcolor{cwlightred}7.14\% \\
    & stat sig negative & \cellcolor{cwnormal}0.00\%   & \cellcolor{cwlightred}3.57\% & \cellcolor{cwlightred}3.57\% \\
\bottomrule
\end{NiceTabular}}
\vspace{1em}
\caption{
Confusion matrix of the statistical conclusions (at significance level $\alpha=0.1$)
drawn using the \textit{Average GMV per User} (Y-axis) and the predicted (X-axis) lLTV Uplift for 28 high-impact A/B tests.}
\label{tbl:sign_validation_avg_gmv}
\end{table}
\end{comment}

В случаях, когда обе метрики фиксируют статистически значимый эффект,
направления выводов совпадают; статистически значимых противоречий
между ними не наблюдается. Дополнительно примерно в $21.4\%$
экспериментов \textbf{Orbitals} обнаруживает значимый эффект там, где
средняя выручка на пользователя его не различает. Эти случаи
соответствуют изменениям, воздействующим не на немедленную
монетизацию, а на структуру переходов между орбиталями (улучшения
пользовательских сценариев, элементы программ лояльности и т.\,п.).


\subsection{Анализ конкретного эксперимента: BNPL}
\label{subsec:bnpl_case}

Запуск опции «Buy Now, Pay Later» позволяет проиллюстрировать
центральный тезис раздела~\ref{sec:framework}: эффект
экспериментального воздействия проявляется как сдвиг распределения
пользователей между орбиталями. На горизонте 41 дня в тестовой
группе наблюдались повышенные вероятности переходов в орбитали
\texttt{active\&paid} и \texttt{whale} относительно контрольной группы
(см.~табл.~\ref{tbl:bnpl_transitions}).

\begin{comment}
\begin{table}[h]
\centering
\scalebox{0.65}{
\begin{NiceTabular}{c l r r r r r}
\toprule
 & & \multicolumn{5}{c}{post\_orbital} \\
\cmidrule{3-7}
\Block{7-1}{\rotate \centering orbital} & & weak, & paid, & active, & active\&paid, & whale, \\
\midrule
 & out / 2.21\%          & \colorcell{-0.0327} & \colorcell{0.6066} & \colorcell{-1.0299} & \colorcell{0.3735} & \colorcell{0.0428} \\
 & weak / 14.50\%        & \colorcell{-1.0884} & \colorcell{0.4685} & \colorcell{-0.7679} & \colorcell{1.5517} & \colorcell{0.0543} \\
 & paid / 4.53\%         & \colorcell{-0.1447} & \colorcell{0.0304} & \colorcell{-1.6038} & \colorcell{1.3848} & \colorcell{0.2169} \\
 & active / 22.66\%      & \colorcell{-0.6261} & \colorcell{0.2539} & \colorcell{-2.1411} & \colorcell{2.3358} & \colorcell{0.2101} \\
 & active\&paid / 50.42\%& \colorcell{-0.2525} & \colorcell{0.0202} & \colorcell{-1.6229} & \colorcell{1.2491} & \colorcell{0.6229} \\
 & whale / 5.68\%        & \colorcell{0.0430} & \colorcell{0.0423} & \colorcell{-0.5891} & \colorcell{-1.2199} & \colorcell{1.7389} \\
\bottomrule
\end{NiceTabular}
}
\caption{Разность матриц переходов (тест минус контроль) для эксперимента BNPL.}
\label{tbl:bnpl_transitions}
\end{table}
\end{comment}

В терминах раздела~\ref{subsec:ab_orbitals} это соответствует
систематическому различию эмпирических матриц переходов
$\widehat{P}^{(B)}$ и $\widehat{P}^{(A)}$, приводящему к различию
итоговых распределений $D_B(T_E)$ и $D_A(T_E)$ и, как следствие, к
ненулевому значению $r_{V\text{-corr}}$. Существенно, что наибольшие
различия наблюдаются у пользователей со средним уровнем трат, что
даёт интерпретируемый продуктовый вывод о потенциале таргетированного
онбординга для этой аудитории.

Frequency‑based вариант методики предсказал
$\hat{r}_{V\text{-corr}} = +1.076\%$; фактическое значение аналогичного
частотного эффекта на горизонте $\mathcal{F}(T_E, 30)$ составило
$+0.789\%$, относительная ошибка прогноза — порядка $36\%$
(см.~рис.~\ref{fig:bnpl_gmv_frequency},
панель~\ref{fig:bnpl_frequency}). Несмотря на завышение абсолютной
величины, метод корректно восстанавливает знак и порядок эффекта при
существенно более узких бутстрэп‑интервалах, чем у прямого эталона.
С практической точки зрения возможность раннего и устойчивого
выявления направления долгосрочного эффекта согласуется с ролью
орбитального оценивателя, обозначенной в~\ref{sec:estimation_tasks}:
повышение чувствительности теста при ограниченном объёме выборки.

В этом сценарии краткосрочная метрика \textit{Average GMV per User}
не обладает достаточной разрешающей способностью, чтобы отделить
первичное принятие новой опции от формирования устойчивого паттерна
поведения. В отличие от неё, $r_{V\text{-corr}}$ непосредственно
отражает изменение траекторий пользователей между орбиталями.

\begin{comment}
\begin{figure}[h]
    \centering
    \begin{subfigure}[t]{0.48\linewidth}
        \includegraphics[width=\textwidth]{data/new/05_bnpl_gmv.png}
        \caption{revenue-based}
        \label{fig:bnpl_gmv}
    \end{subfigure}
    \hfill
    \begin{subfigure}[t]{0.48\linewidth}
        \includegraphics[width=\textwidth]{data/new/05_bnpl_frequency.png}
        \caption{frequency-based}
        \label{fig:bnpl_frequency}
    \end{subfigure}
    \caption{
        Histogram plots of \textit{bootstrap} samples with
        observed (Y-axis) and predicted (X-axis) lLTV Uplift statistics
        for BNPL experiment. The solid line denotes \(y = x\).
    }
    \label{fig:bnpl_gmv_frequency}
\end{figure}
\end{comment}


\subsection{Сопоставление с эталонными методами}
\label{subsec:benchmark_comp}

В этом разделе закрывается обещание, сформулированное в
разделе~\ref{sec:estimation_tasks}: компромисс смещение–дисперсия
квантифицируется на реальных данных. Сопоставляются три семейства
оценок:
\begin{itemize}
    \item орбитальные оцениватели $r_V$ и $r_{V\text{-corr}}$;
    \item прямой эталон по реализованному $R_u(T_E)$
          (раздел~\ref{sec:estimation_tasks}, пункт~(1));
    \item сильная индивидуальная предсказательная модель
          (раздел~\ref{sec:estimation_tasks}, пункт~(2)).
\end{itemize}


\paragraph{Сильная предсказательная модель.}

В качестве индивидуальной модели используется
\texttt{CatBoostRegressor}, обучаемый предсказывать $R_u(T_E)$ по тому
же набору поведенческих признаков, что и орбитальная методика
(раздел~\ref{subsubsec:feature_engineering}). Обучение проводится на
выборке порядка сотен тысяч пользователей с разбиением $70/30$;
функция потерь и метрика — RMSE; число итераций бустинга — $100$,
скорость обучения — $0.2$, глубина деревьев — $6$. Для эксперимента
модель применяется к пользователям обеих групп в момент $T_E$,
порождая предсказания $\tilde{R}_u$, по которым формируется
относительная разница $r_{\text{pred}}$ и её скорректированный аналог
$r_{\text{pred-corr}}$ (по аналогии с $r_{V\text{-corr}}$).

\begin{comment}
\begin{equation*} \label{eq:pred_rel_diff}
    r_{\text{pred}}(B \parallel A)
    = \frac{\tilde{V}(T_E; B) - \tilde{V}(T_E; A)}{\tilde{V}(T_E; A)}.
\end{equation*}
\end{comment}


\paragraph{Сценарии и результаты.}

Сравнение проводится на двух сценариях:
\begin{description}
    \item[A/A‑тест:] нулевой истинный эффект; объём выборки сопоставим
          с реальным экспериментом. Позволяет оценить уровень ложных
          срабатываний и дисперсию оценивателей.
    \item[A/B‑эксперимент BNPL:] ненулевой эффект; объём — десятки
          тысяч пользователей в каждой группе. Позволяет сопоставить
          чувствительность методов.
\end{description}

Результаты (см.~рис.~\ref{fig:histrogram}) воспроизводят
теоретический прогноз раздела~\ref{sec:estimation_tasks}:
\begin{itemize}
    \item в A/A‑сценарии орбитальные оцениватели остаются близкими
          к нулю и обладают существенно меньшей дисперсией, чем
          прямой эталон и индивидуальная модель; в этих
          альтернативах фиксируются ложные положительные
          бутстрэп‑эффекты, ожидаемые в силу высокой дисперсии и
          модельных смещений соответственно;
    \item в BNPL‑сценарии $r_{V\text{-corr}}$ согласуется с
          $r_{\text{pred-corr}}$ по знаку и порядку величины эффекта,
          но обладает существенно меньшей дисперсией, чем оба
          эталона.
\end{itemize}

Тем самым эмпирически воспроизведён компромисс, заявленный в
разделе~\ref{sec:estimation_tasks}: небольшое модельное смещение
(в случае приближённой выполнимости мультипликативного допущения
из~\ref{subsec:lltv}) обменивается на существенное снижение дисперсии
за счёт перехода от индивидуальных предсказаний $R_u$ к
низкоразмерному параметру $\{V_c\}_{c\in\mathcal{C}}$.

С целью воспроизводимости результатов открытая реализация
\textbf{Orbitals} опубликована вместе с обезличенным набором дневных
пользовательских логов более чем 1~млн пользователей. Код и данные
доступны в публичном
репозитории.\footnote{\href{https://github.com/orbitals-acc/orbitals-framework}{https://github.com/orbitals-acc/orbitals-framework}}


\subsection{Обсуждение и ограничения}
\label{subsec:limitations}

Эмпирические результаты согласуются с теоретическими ожиданиями
предыдущих глав по трём ключевым пунктам: подтверждается
мультипликативное допущение из раздела~\ref{subsec:lltv};
орбитальный оцениватель снижает дисперсию относительно прямого
эталона и индивидуальной модели, как заявлено в
разделе~\ref{sec:estimation_tasks}; устойчивость сегментации,
постулированная в разделах~\ref{subsec:orbitals_def}
и~\ref{subsubsec:model_training}, не нарушается на рассматриваемом
горизонте.

Подход опирается на два явных допущения, ограничивающих область его
применимости:
\begin{description}
    \item[Стабильность сегментации
          (\ref{subsec:orbitals_def},~\ref{subsubsec:model_training}):]
          поведенческие определения орбиталей сохраняют сопоставимый
          смысл во времени. Выраженная сезонность, изменение состава
          аудитории или крупные изменения продукта могут потребовать
          переобучения ансамбля и пересмотра $\mathcal{C}$ и
          $\{\hat{V}_c\}$.
    \item[Симметрия внешних факторов:] внешние воздействия в окне
          $\mathcal{F}(T_E, T_F)$ влияют на контрольную и тестовую
          группы симметрично. При нарушении этого условия
          $r_{V\text{-corr}}$ может оказаться смещённым.
\end{description}

Отдельной осторожности требует интерпретация экспериментов,
целенаправленно воздействующих на узкие сегменты пользовательской
базы (например, на отдельные орбитали): в таких случаях
мультипликативное допущение может локально нарушаться и требовать
дополнительной диагностики.

Несмотря на отмеченные ограничения, проведённое исследование позволяет
заключить, что при корректной калибровке \textbf{Orbitals} реализует
анонсированный в~\ref{sec:estimation_tasks} компромисс между
смещением и дисперсией и пригоден в качестве интерпретируемого
инструмента ранней оценки долгосрочного эффекта продуктовых
интервенций.